\section{An exact nonlinear flux law and compatible radial source equations}
\label{sec:r12-radial-flux-source-inversion}

Keep the original $n$-polydisc, $n>q\ge1$, and the single
fixed original background $S$ of \textup{(PO3)}.  This is
an explicit family of allowed backgrounds in every remaining
degree, not a replacement of the arbitrary background in the
general theorems.  Write $d=q+1$ and retain the original
radial parameter $t$ and diagonal Gauss map $\sigma$.

For a positive Radon measure $\nu$ on $(0,\infty)$ with
finite tails, use the \emph{open-tail} convention
\[
 H_\nu(t)=\nu((t,\infty)),\qquad t>0.
 \tag{RF1}
\]
Thus $H_\nu$ is right continuous and $-dH_\nu=\nu$
as distributions on $(0,\infty)$.  An atom at $t$ is
excluded from $H_\nu(t)$, but included in its left limit.
This convention is important in the clipping formulas below.

\begin{definition}[Admissible radial slopes and source fluxes]
\label{def:r12-radial-admissible-flux}
Let $\mathscr A_q$ be the cone of finite-valued,
nonnegative, nonincreasing right-continuous functions
$A:(0,\infty)\to[0,\infty)$ such that
\[
                 \int_0^\infty e^{-t}A(t)^d\,dt<\infty.
 \tag{RF2}
\]
Set
\[
 F_A(t)=\int_0^t A(s)\,ds,\qquad
 u_A(x)=-F_A(t(x))\quad(x\ne o),\qquad
 u_A(o)=-F_A(\infty).
 \tag{RF3}
\]
The value at $o$ may be $-\infty$.  A source $\nu$ is
\emph{radially flux-compatible} if it has finite tails,
the function $e^tH_\nu(t)$ is nonincreasing, and
\[
              \int_0^\infty e^{t/q}H_\nu(t)^{d/q}\,dt<\infty.
 \tag{RF4}
\]
These are conditions on the ordinary one-dimensional
Radon datum.  They do not invoke an unknown potential
or an unspecified solution of a source equation.
\end{definition}

\begin{theorem}[Exact radial flux and finite-energy inversion]
\label{thm:r12-exact-radial-flux-inverse}
Every $A\in\mathscr A_q$ defines an actual extended
classical psh potential $u_A$ in the original no-loss
zero-boundary first-energy class.  Its original top source is
\[
 \mu_{u_A}=\sigma_*\nu_A,\qquad
 H_{\nu_A}(t)=e^{-t}A(t)^q.
 \tag{RF5}
\]
Its energy and every mixed first energy have the exact formulas
\[
 \begin{split}
 e(u_A)&=\int_0^\infty e^{-t}A(t)^d\,dt,\\
 B(u_{A_0},\ldots,u_{A_q})
     &=\int_0^\infty e^{-t}\prod_{i=0}^q A_i(t)\,dt.
 \end{split}
 \tag{RF6}
\]
All source identities hold on every original compact smooth
scalar test and every actual moment refinement.

The source map from these radial potentials is a bijection
onto the flux-compatible Radon data.  The inverse is explicit:
\[
 A_\nu(t)=\bigl(e^tH_\nu(t)\bigr)^{1/q},\qquad
 u_\nu=-\int_0^{t(\cdot)}
                     \bigl(e^sH_\nu(s)\bigr)^{1/q}ds.
 \tag{RF7}
\]
For such a datum $u_\nu$ is the unique solution of
$\widehat\mu_v=\sigma_*\nu$ in the \emph{full} strong
classical-approximation cone $\mathscr C$, not just among
radial candidates.  The flux condition is asserted necessary
only within the displayed radial class; no classification of
all arbitrary nonradial source measures is claimed.
\end{theorem}

\noindent\textbf{Status: \texttt{NEW\_PROVED}} for the
explicit fixed allowed background.  The arbitrary-background
bounded calculus, no-loss extension, first-energy mixed
calculus, and original radial-shell calibration are frozen
\texttt{INTERNAL\_INPUT}.

\begin{proof}
On every interval $(0,R)$ the exponential weight is bounded
below, so Holder's inequality makes $F_A(R)$ finite.
The function $F_A$ is nonnegative, increasing and concave,
with $F_A(0)=0$.  The scalar function
$z\mapsto-F_A(-z)$ on $z<0$ is convex and nondecreasing.
The frozen radial/convex composition calculus therefore gives
the actual extended classical psh function in \textup{(RF3)}.
One can also construct it by decreasing the bounded classical
radial potentials with slopes
\[
                  A_{N,R}(t)=\min(A(t),N)\mathbf1_{(0,R)}(t).
 \tag{RF8}
\]
Each of these extends across the boundary by a finite linear
continuation of its radial scalar function.  The limits are
taken on the original space, with no extension of $K$.
Their finite values at every finite positive depth show
that the limit is not identically $-\infty$.  The boundary value of $u_A$ is zero
because $F_A(t)\to0$ as $t\downarrow0$, and its strict
negative sublevels are compactly contained smaller polydiscs.

For the moment take a bounded radial slope, clipped to zero
after a finite depth.  Polarization of the same original
radial-shell calibration used in \textup{(PO6)} gives
\[
 H_{T_{[V],n}(u_A^{[q]},\psi_s^{[n-q]})}(t)
                      =A(t)^q\mathbf1_{(0,s)}(t)
 \tag{RF9}
\]
at all continuity points.  Here the tail is taken along the
original diagonal Gauss ray; the formula asserts the full
original measure identity.  Indeed the top measure of
$a u_A+b\psi_s$ has tail
$(aA(t)+b\mathbf1_{(0,s)}(t))^n$ away from jump points.
The coefficient of $a^qb^{n-q}$ has the same binomial
factor on both sides of mixed multilinearity, so there is
no additional factorial.  Right continuity determines the
measure and all jump masses uniquely.

Insert the fixed background \textup{(PO3)}.  The coefficient
Fubini argument of \textup{(PO5)}, applied on each fixed
compact original stage, gives
\[
          H_{\mu_{u_A}}(t)
                =A(t)^q\int_t^\infty e^{-s}ds
                =e^{-t}A(t)^q.
 \tag{RF10}
\]
For a slope not bounded near zero, first fix $R$.  One has
\[
 \sup_V\left|u_{A_{N,R}}-u_{A\mathbf1_{(0,R)}}\right|
        \le\int_0^R(A-\min(A,N))\,dt\longrightarrow0.
\]
Frozen bounded uniform stability therefore gives the
clipped original source, including on compact sets
containing the origin.  Away from clipping depth $R$
its tail is the stated one.  Local boundedness on
every $(\varepsilon,R)$ gives the same formula there.
Only after this fixed-$R$ passage do we let $R$ tend
to infinity, using \textup{(RF11)} below to control
the entire collapsed tail and the value-truncation loss.
A compact subset of $V$ stays a positive radial depth
away from the outer boundary, so the zero-depth boundary
is not silently added to $V$.

Let $H_A(t)=e^{-t}A(t)^q$.  It is right continuous,
nonincreasing, finite for $t>0$, and tends to zero at
infinity.  It therefore defines a unique positive Radon
measure $\nu_A=-dH_A$ with the displayed open tail.
Clipping the slope at depth $R$ collapses its entire
remaining source to that depth:
\[
 \nu_{A\mathbf1_{(0,R)}}
       =\mathbf1_{(0,R)}\nu_A+\nu_A([R,\infty))\delta_R,
 \qquad
 \nu_A([R,\infty))=e^{-R}A(R-)^q.
 \tag{RF11}
\]
There is no ambiguity when $A$ jumps at $R$: the left
limit includes the pre-existing atom.  For $R\ge1$,
this collapsed mass is bounded by $e^{-R}A(1-)^q$
after choosing any fixed earlier positive depth instead
of $1$ if necessary.  In particular it tends to zero.

If $F_A$ is unbounded, its value truncation at $-j$
corresponds to the unique depth $R_j$ with $F_A(R_j)=j$;
on any possible plateau one uses the first such depth,
which gives the same clipped potential.  When $F_A$
is unbounded its slope cannot vanish at a finite depth,
so these depths tend to infinity.  Formula \textup{(RF11)}
proves absence of truncation loss.  If $F_A$ is bounded,
the potential is already bounded and there is no pole
loss.  Thus the frozen no-loss class and original source
in \textup{(RF5)} are established without inferring
membership from a formally finite nonpluripolar energy.

Integration by parts now gives
\[
              \int_0^\infty F_A\,d\nu_A
                   =\int_0^\infty A(t)H_A(t)\,dt.
 \tag{RF12}
\]
For completeness, both endpoint terms vanish.  At infinity
$F_A$ grows at most linearly and $H_A$ is exponentially
bounded.  At zero, monotonicity and Holder imply
\[
 F_A(t)A(t)^q
 \le \left(t^{q/d}\Bigl(\int_0^t A^d\Bigr)^{1/d}\right)
       \left(t^{-q/d}\Bigl(\int_0^t A^d\Bigr)^{q/d}\right)
 =\int_0^t A^d\longrightarrow0.
\]
The exponential factor does not affect this limit.
Substitution in \textup{(RF12)} proves the first identity
in \textup{(RF6)}, hence actual first-energy membership.
The second follows by polarizing that identity in
$A_0,\ldots,A_q$, using the frozen symmetric mixed first
energy.  All integrals are finite by ordinary Holder.

Conversely, the flux hypotheses make $A_\nu$ a member
of $\mathscr A_q$: the integral in \textup{(RF4)} is
exactly \textup{(RF2)} after substitution.  Formula
\textup{(RF5)} constructs $u_\nu$ with the specified
original Radon source.  The right-continuous tail recovers
$A$, and integration with boundary value zero recovers
the radial potential, proving the radial bijection.

Finally suppose any $v\in\mathscr C$ has the same
original source measure as $u_\nu$.  The literal mixed
energy extension identifies their curvature functionals
on every cone member, hence on $\mathbb D_S$.
The strong discrepancy is zero and frozen separation
gives $v=u_\nu$ as relative classes.  This proves
uniqueness in the full cone without asserting radiality
of arbitrary other source data.
\end{proof}

\begin{corollary}[A linear compatibility condition for finite radial data]
\label{cor:r12-finite-radial-data-compatibility}
Let $\nu$ be a finite positive Radon measure on $(0,\infty)$,
of total mass $M$.  It is flux-compatible if and only if
\[
                         \nu\ge H_\nu(t)\,dt.
 \tag{RF13}
\]
Equivalently, $e^tH_\nu(t)$ is nonincreasing; no further
finite-energy condition is needed in this finite-mass case.
The solution \textup{(RF7)} has bounded initial slope
$A_\nu(0+)=M^{1/q}$, belongs to every finite positive
power-energy class, and satisfies
\[
             e_\alpha(u_\nu)
                 \le M^{1+\alpha/q}\Gamma(\alpha+1)
                 \qquad(\alpha>0).
 \tag{RF14}
\]
The same bound holds for its bounded radial truncations.
\end{corollary}

\noindent\textbf{Status: \texttt{DIRECT\_COROLLARY}} of
the exact source theorem, with an explicit datum criterion.

\begin{proof}
Distributional differentiation gives
\[
                 d(e^tH_\nu)=e^t(H_\nu\,dt-\nu).
\]
Thus \textup{(RF13)} is exactly monotonicity of the
right-continuous flux.  Since $H_\nu(0+)=M$,
compatibility gives $H_\nu(t)\le Me^{-t}$ and
$A_\nu(t)\le M^{1/q}$.  This proves the energy
integrability and the initial-slope assertion.  The
bound $F_{A_\nu}(t)\le M^{1/q}t$ and layer cake yield
\[
 \int F_{A_\nu}(t)^\alpha\,d\nu(t)
 \le M^{\alpha/q}\alpha\int_0^\infty
                              t^{\alpha-1}H_\nu(t)dt
 \le M^{1+\alpha/q}\Gamma(\alpha+1).
\]
The no-loss proof is already part of the theorem.
In \textup{(RF11)} each source tail is moved to the
clipping depth and its potential value becomes the
corresponding smaller clipped value.  Its power
energy is therefore no larger than the original
integral, proving the last assertion.
\end{proof}

\begin{proposition}[Exact discrepancy and stability of the inverse]
\label{prop:r12-radial-flux-stability}
For $A,B\in\mathscr A_q$ one has
\[
 I(u_A,u_B)=\int_0^\infty
               e^{-t}(A-B)(A^q-B^q)\,dt.
 \tag{RF15}
\]
With all norms below taken in $L^d(e^{-t}dt)$,
\[
 \|A-B\|_d^d\le I(u_A,u_B)
   \le q(\|A\|_d+\|B\|_d)^{q-1}\|A-B\|_d^2.
 \tag{RF16}
\]
Consequently strong convergence in this radial class
is exactly convergence of slopes in this weighted
$L^d$ space.  For compatible sources it is equivalently
\[
 \int_0^\infty e^{t/q}
       |H_{\nu_j}(t)^{1/q}-H_\nu(t)^{1/q}|^d\,dt
                                                        \longrightarrow0.
 \tag{RF17}
\]
In particular, vaguely convergent compatible finite
data of uniformly bounded total mass have strongly
convergent solutions, provided the vague limit is
taken on $(0,\infty)$; the conclusion includes loss
of source mass through the excluded outer boundary.
\end{proposition}

\noindent\textbf{Status: \texttt{NEW\_PROVED}} for
the exact radial nonlinear inverse; no general data
continuity outside this fixed-background family is inferred.

\begin{proof}
Apply the original discrepancy identity and then
integrate by parts, with the vanishing endpoint
terms justified as in \textup{(RF12)}.  This gives
\[
 I(u_A,u_B)
  =\int(F_A-F_B)d(\nu_A-\nu_B)
  =\int(A-B)(H_{\nu_A}-H_{\nu_B})dt,
\]
which is \textup{(RF15)}.  For nonnegative scalars
$a,b$ and integer $q\ge1$,
\[
 |a-b|^{q+1}\le(a-b)(a^q-b^q)
                      \le q(a+b)^{q-1}|a-b|^2.
\]
Integrate the first inequality.  Holder with exponents
$d/(q-1)$ and $d/2$, followed by Minkowski, proves
the second; for $q=1$ it is simply an equality.
The energy identity in \textup{(RF6)} makes the
topological equivalence immediate.  Substitution
of \textup{(RF7)} gives \textup{(RF17)}.

For the final statement let the total masses be
bounded by $M$.  Compatibility gives the uniform
tail bound $H_{\nu_j}(t)\le Me^{-t}$, so no mass
can escape to infinity.  Vague convergence therefore
implies tail convergence at every continuity point
of the limit tail, by first restricting to a compact
interval and then using the exponential bound.
The limit is finite and compatible: the distributional
inequality \textup{(RF13)} passes locally to the limit,
and its tails obey the same bound.  Almost-everywhere
convergence of $A_{\nu_j}$, their uniform bound
$M^{1/q}$, and dominated convergence in $L^d(e^{-t}dt)$
prove strong convergence.  Source mass may escape to
zero, which is outside $V$'s ordinary radial domain;
the proof does not assert preservation of total mass
under this vague convergence.
\end{proof}

\subsection{An explicit nonpolar datum with no full-cone source solution}

\begin{theorem}[A positive-capacity atom and its exact variational replacement]
\label{thm:r12-positive-capacity-atom-source-obstruction}
Fix $R>0$, write $E_R=1-e^{-R}$, and let
\[
 \lambda_R=\delta_{\sigma(R)},\qquad
 \ell_R(h)=-h(\sigma(R)),\qquad
 c_R=\left(\frac R{E_R}\right)^{1/q},\qquad
 v_R=c_R\psi_R.
 \tag{RF18}
\]
The datum is a literal positive Radon measure supported
at a point of positive full outer capacity.  It is
dominated by the actual bounded source
$e^R\mu_{\psi_R}$, so $\ell_R\in\mathbb D_{S,\mathrm s}^*$.
Its exact norm and its unique variational maximizer
over the entire completed cone are
\[
       \|\ell_R\|_* =\frac R{E_R^{1/d}},\qquad
       \mathop{\rm argmax}_{a\in\mathscr C}
                  \{\ell_R(a)-e(a)/d\}=v_R.
 \tag{RF19}
\]
Nevertheless there is no $a\in\mathscr C$ satisfying
$\widehat\mu_a=\lambda_R$ on all original tests.
Indeed the maximizer has the different exact source
\[
 \widehat\mu_{v_R}
     =\frac R{E_R}\,
           \sigma_*\bigl(\mathbf1_{(0,R)}e^{-s}ds+e^{-R}\delta_R\bigr).
 \tag{RF20}
\]
Thus positive-capacity support, domination by an actual
bounded source, and all finite power-energy bounds on
the variational solution do not suffice for the
unrestricted source equation.
\end{theorem}

\noindent\textbf{Status: \texttt{NEW\_PROVED}}.  The
unrestricted source-existence claim under only the
listed conditions is \texttt{FALSE}, in every $q\ge1$.
The background remains the single explicit allowed
current \textup{(PO3)}.

\begin{proof}
Positivity of capacity and the literal atom bound
follow from \textup{(PO4)}.  Frozen strong tester
continuity against the fixed actual source proves
strong continuity of $\ell_R$; this is literal
evaluation at a non-null point, not the nonliteral
null-fiber evaluation considered elsewhere.

For arbitrary $a\in\mathscr C$ put
$F_a(t)=-a(\sigma(t))$.  The completed trace theorem
makes $F_a$ nonnegative and concave with $F_a(0)=0$.
Its chord bound on $[0,R]$ gives
\[
 \begin{split}
 B(a,\psi_R^{[q]})
   &=\int_0^R F_a(s)e^{-s}ds+e^{-R}F_a(R)\\
   &\ge\frac{F_a(R)}R
             \left(\int_0^R se^{-s}ds+Re^{-R}\right)
     =\frac{E_R}{R}F_a(R).
 \end{split}
 \tag{RF21}
\]
Mixed Holder bounds the left side by $r(a)E_R^{q/d}$.
Consequently $\ell_R(a)\le R E_R^{-1/d}r(a)$,
with equality at a positive multiple of $\psi_R$.
The cone/difference dual norm identity proves the
first assertion of \textup{(RF19)}.  Optimizing
$cr\mapsto\ell_R(c\psi_R)-e(c\psi_R)/d
=cR-c^dE_R/d$ gives $c=c_R$, and the same norm
bound proves this is a global maximizer on the full
cone.  Frozen strict midpoint energy separation
proves uniqueness.  No radiality assumption is
imposed on the competitors in this argument.

Homogeneity of the original source gives
\textup{(RF20)}.  Its diffuse positive measure on
every compact radial subinterval of $(0,R)$ makes
it different from $\lambda_R$.  This difference is
visible on original compact smooth scalar tests:
choose a nonnegative smooth compact real function
$\chi$ supported in $(-R,0)$ and put, on the
coordinate-unit torus,
$\eta=\prod_{i=1}^n\chi(\log|T_i|)$, extended
by zero near all coordinate zero loci.  Its support
is an ordinary compact polyannulus in $V$,
$\eta(\sigma(R))=0$, and its integral against
\textup{(RF20)} is strictly positive when $\chi$
is not identically zero.  Thus the discrepancy is
not merely a failure on auxiliary or delta tests.

If an $a\in\mathscr C$ solved the source equation,
then $J(a)=\ell_R$ by literal integration on the
energy-difference space.  The differentiable energy
support inequality makes such an $a$ a variational
maximizer.  Uniqueness would force $a=v_R$, in
contradiction with \textup{(RF20)} and the preceding
original test.  This proves full-cone nonexistence.
\end{proof}

\begin{corollary}[The normal residual is exactly a concavity defect]
\label{cor:r12-radial-atom-normal-concavity}
For the datum and maximizer in \textup{(RF18)}, the
normal functional $n_R=\ell_R-J(v_R)$ satisfies,
for every $a\in\mathscr C$,
\[
 n_R(a)=-\frac R{E_R}\int_0^R
          \left(F_a(s)-\frac{s}{R}F_a(R)\right)e^{-s}ds\le0.
 \tag{RF22}
\]
It vanishes precisely on those cone members whose
actual radial trace is linear on $[0,R]$.
In particular $n_R(v_R)=0$, but $n_R$ is nonzero
as an original-test functional.
\end{corollary}

\noindent\textbf{Status: \texttt{NEW\_PROVED}} exact
geometric description of one full-cone normal face.

\begin{proof}
Use \textup{(PO4)} and $c_R^q=R/E_R$ to expand
$J(v_R)(a)$.  Since
$\int_0^Rse^{-s}ds=E_R-Re^{-R}$, the result is
\textup{(RF22)}.  The integrand is nonnegative by
concavity and continuous, so its weighted integral
vanishes exactly when it vanishes identically.
That is precisely the asserted linearity.  The
original test constructed above gives a nonzero
value of the residual, so its negativity on the
potential cone is not an equality on the original
smooth test sheaf.
\end{proof}

\paragraph{What this advances, and what it does not.}
The flux condition gives an explicit compatible
nonlinear source class, a literal actual solution,
full-cone uniqueness, and a quantified strong inverse
in every remaining degree.  The atom computation
identifies a nonzero normal residual as a concrete
concavity defect, rather than assuming it vanishes
because the datum is positive or capacity controlled.
It is a full-dimensional stress test for the general
source problem, not a universal classification of
arbitrary-$S$ sources and not a proof of OPEN-GQ or
OPEN-CTO.  No auxiliary analytic field, change of
absolute value, or restriction to toric test forms
is used in any displayed current equality.
